\documentclass[11pt,a4paper]{article}

% Template visual Matriz Aprova. Compile com XeLaTeX ou LuaLaTeX.
\usepackage{fontspec}
\usepackage{polyglossia}
\setmainlanguage{portuguese}
\IfFontExistsTF{Aptos}{
  \setmainfont{Aptos}
  \setsansfont{Aptos}
}{
  \IfFontExistsTF{Calibri}{
    \setmainfont{Calibri}
    \setsansfont{Calibri}
  }{
    \IfFontExistsTF{TeX Gyre Heros}{
      \setmainfont{TeX Gyre Heros}
      \setsansfont{TeX Gyre Heros}
    }{
      \setmainfont{Latin Modern Sans}
      \setsansfont{Latin Modern Sans}
    }
  }
}
\usepackage[a4paper,margin=2cm,headheight=16pt]{geometry}
\usepackage{graphicx}
\usepackage{xcolor}
\usepackage{tikz}
\usepackage{tcolorbox}
\usepackage{enumitem}
\usepackage{booktabs}
\usepackage{tabularx}
\usepackage{amssymb}
\usepackage{fancyhdr}
\usepackage{titlesec}
\usepackage{hyperref}

\definecolor{matrizlime}{HTML}{CBFF4D}
\definecolor{matrizink}{HTML}{101313}
\definecolor{matrizmuted}{HTML}{596363}
\definecolor{matrizline}{HTML}{DDE3DF}
\definecolor{matrizsoft}{HTML}{F3F7EC}

\hypersetup{colorlinks=true,linkcolor=matrizink,urlcolor=matrizink}
\pagestyle{fancy}
\fancyhf{}
\lhead{\textsf{\small MATRIZ APROVA}}
\rhead{\textsf{\small APOSTILA}}
\cfoot{\textsf{\small \thepage}}
\renewcommand{\headrulewidth}{0.4pt}
\renewcommand{\headrule}{\hbox to\headwidth{\color{matrizline}\leaders\hrule height \headrulewidth\hfill}}

\titleformat{\section}{\Large\bfseries\color{matrizink}}{\thesection}{0.6em}{}
\titleformat{\subsection}{\large\bfseries\color{matrizink}}{\thesubsection}{0.6em}{}
\setlist{nosep,leftmargin=*}

\newtcolorbox{foco}{colback=matrizsoft,colframe=matrizlime,boxrule=1pt,arc=3pt,left=10pt,right=10pt,top=8pt,bottom=8pt}
\newtcolorbox{lei}{colback=matrizink,colframe=matrizink,coltext=white,boxrule=0pt,arc=3pt,left=10pt,right=10pt,top=8pt,bottom=8pt}
\newtcolorbox{alerta}{colback=white,colframe=matrizline,boxrule=0.6pt,arc=3pt,left=10pt,right=10pt,top=8pt,bottom=8pt}

% Logo usada na capa. Caminho relativo à pasta onde o .tex é compilado
% (materiais/<disciplina>/). Ajuste se a estrutura mudar.
\newcommand{\logomatriz}{../../public/logo.png}

\begin{document}

\begin{titlepage}
  \thispagestyle{empty}
  \begin{center}
    \includegraphics[width=0.55\linewidth]{\logomatriz}\par
  \end{center}
  \vspace{2.2cm}
  {\sffamily\fontsize{28}{32}\selectfont\bfseries Apostila — Direito do Trabalho: Relação de emprego e requisitos do vínculo\par}
  \vspace{0.4cm}
  {\sffamily\large TRT e OAB\par}
  \vfill
  \begin{foco}
    \textsf{\textbf{Foco da apostila}}\\[3pt]
    não basta existir pagamento. O vínculo depende da análise conjunta dos requisitos e da realidade da prestação, não apenas do nome dado ao contrato.
  \end{foco}
  \vspace{0.7cm}
  {\sffamily\small Atualizado em 16 de agosto de 2026\quad ·\quad FCC e FGV\par}
  \vspace{1cm}
  {\sffamily\small matrizaprova.com\par}
\end{titlepage}

\tableofcontents
\newpage

% Conteúdo gerado entra aqui.
\section*{O que cai}

O vínculo de emprego é reconhecido quando estão presentes, conjuntamente, os requisitos jurídicos da relação empregatícia. As questões costumam explorar:

\begin{enumerate}
  \item pessoa física;
  \item pessoalidade;
  \item não eventualidade ou habitualidade;
  \item onerosidade;
  \item subordinação jurídica;
  \item diferença entre empregado, empregador, autônomo, eventual e estagiário.
\end{enumerate}

\textbf{Regra de prova:} não basta existir pagamento. O vínculo depende da análise conjunta dos requisitos e da realidade da prestação, não apenas do nome dado ao contrato.

\section*{Teoria essencial}

\subsection*{1. Empregado}

O art. 3º da CLT considera empregado toda pessoa física que prestar serviços de natureza não eventual a empregador, sob a dependência deste e mediante salário.

A leitura do dispositivo permite identificar os requisitos clássicos:

\begin{itemize}
  \item pessoa física;
  \item pessoalidade;
  \item não eventualidade;
  \item onerosidade;
  \item subordinação jurídica.
\end{itemize}

A ausência de qualquer elemento pode afastar a relação de emprego, embora seja necessário analisar o caso concreto e a legislação específica aplicável.

\subsection*{2. Pessoa física}

Empregado é pessoa natural. Uma pessoa jurídica pode prestar serviço a outra empresa por contrato civil ou empresarial, mas não é empregada como pessoa jurídica.

O fato de o trabalhador possuir CNPJ ou emitir nota fiscal não elimina, por si só, a possibilidade de reconhecimento do vínculo. Se a prestação real possuir os requisitos dos arts. 2º e 3º da CLT, a forma documental não deve ocultar a realidade.

\subsection*{3. Pessoalidade}

O trabalho deve ser prestado pela própria pessoa contratada. O empregado não pode se substituir livremente por terceiro sem autorização do empregador.

A pessoalidade não exige que o trabalhador nunca possa faltar. Férias, licenças e substituições autorizadas não eliminam o requisito. O ponto central é a impossibilidade de substituição livre e permanente por escolha do prestador.

\subsection*{4. Não eventualidade}

Não eventualidade significa inserção do trabalho na dinâmica normal da atividade e expectativa de continuidade. Não é sinônimo obrigatório de trabalho diário.

Um trabalhador pode trabalhar em determinados dias da semana e ainda assim ser empregado, desde que exista habitualidade e continuidade compatíveis com a atividade.

O trabalho eventual, por outro lado, é prestado de forma ocasional, sem integração estável à atividade do tomador. A análise não depende apenas do número de dias, mas da natureza e da continuidade do serviço.

\subsection*{5. Onerosidade}

Onerosidade é a existência de contraprestação econômica pelo trabalho. O salário pode ser pago em dinheiro ou, dentro dos limites legais, envolver utilidades.

Não há vínculo de emprego quando a atividade é verdadeiramente gratuita, como em determinadas atividades voluntárias regidas por legislação própria. Se existir contraprestação disfarçada, a realidade deve ser investigada.

\subsection*{6. Subordinação jurídica}

Subordinação jurídica é a sujeição do trabalhador ao poder de direção do empregador. Ela pode envolver ordens, fiscalização, controle de horários, metas, procedimentos e sanções.

Não é necessário que o empregador dê ordens a todo momento. A subordinação pode ser estrutural: o trabalhador está inserido na organização e nos meios de produção do tomador, submetendo-se à sua dinâmica empresarial.

Subordinação econômica não é requisito autônomo. Um trabalhador pode depender financeiramente de uma empresa sem ser empregado, e pode existir relação de emprego mesmo quando o trabalhador possui outras fontes de renda.

\subsection*{7. Empregador}

O art. 2º da CLT considera empregador a empresa, individual ou coletiva, que, assumindo os riscos da atividade econômica, admite, assalaria e dirige a prestação pessoal de serviço.

O empregador assume o risco do empreendimento. Em regra, não pode transferir ao empregado os riscos econômicos próprios da atividade empresarial.

O empregador exerce poder diretivo, que inclui organizar, fiscalizar e disciplinar o trabalho, sempre dentro da Constituição, da lei, do contrato e dos direitos fundamentais.

\subsection*{8. Trabalho presencial e a distância}

O art. 6º da CLT estabelece que não há distinção entre o trabalho realizado no estabelecimento do empregador, no domicílio do empregado ou a distância, desde que presentes os pressupostos da relação de emprego.

O local de trabalho, isoladamente, não decide se existe vínculo. O trabalho remoto pode ser empregatício; da mesma forma, trabalho presencial pode ser autônomo se não estiverem presentes os requisitos do vínculo.

\subsection*{9. Primazia da realidade}

Na análise trabalhista, os fatos efetivamente praticados têm grande relevância. O contrato pode chamar o trabalhador de parceiro, associado, prestador ou autônomo, mas a nomenclatura não prevalece se a execução demonstrar relação de emprego.

O art. 9º da CLT considera nulos os atos praticados com o objetivo de desvirtuar, impedir ou fraudar a aplicação dos preceitos trabalhistas.

Isso não significa que todo contrato civil ou empresarial seja fraude. É necessário verificar se a autonomia existe de fato e se os requisitos do vínculo estão presentes no cotidiano.

\subsection*{10. Distinções importantes}

\subsubsection*{10.1. Autônomo}

O autônomo organiza sua própria atividade e não se submete à subordinação jurídica típica. Pode prestar serviços com habitualidade, receber pagamento e ser pessoa física; o elemento decisivo é a autonomia real.

\subsubsection*{10.2. Eventual}

O eventual presta serviço ocasional, sem continuidade ou integração permanente à atividade do tomador. Pode receber pagamento, mas falta a não eventualidade.

\subsubsection*{10.3. Estagiário}

O estágio regular possui finalidade educacional e deve observar os requisitos da Lei nº 11.788/2008. O descumprimento das exigências legais pode caracterizar desvirtuamento e gerar consequências trabalhistas.

\subsubsection*{10.4. Voluntário}

O serviço voluntário é atividade não remunerada prestada a entidade pública ou instituição privada sem fins lucrativos, conforme legislação própria. Uma relação apresentada como voluntária, mas com remuneração e requisitos de emprego, deve ser examinada pela realidade.

\section*{Como a banca cobra}

\begin{itemize}
  \item Pessoa jurídica não é empregada, mas o CNPJ não impede vínculo se houver fraude e prestação pessoal subordinada.
  \item Não eventualidade não exige trabalho todos os dias.
  \item Pessoalidade não impede férias ou substituição autorizada.
  \item O local de trabalho não define sozinho a existência do vínculo.
  \item Subordinação jurídica é mais importante que dependência econômica.
  \item O empregador assume os riscos da atividade econômica.
  \item A primazia da realidade pode afastar o rótulo contratual.
  \item Contrato autônomo válido exige autonomia efetiva, não apenas documento.
\end{itemize}

\section*{Exemplos resolvidos}

\subsection*{Exemplo 1}

Uma empresa contrata uma pessoa como "prestadora autônoma", mas exige horário, controla sua rotina, impede substituição e paga mensalmente. Há possibilidade de vínculo?

\textbf{Sim.} O nome do contrato não é decisivo. A prestação pessoal, habitual, onerosa e subordinada pode preencher os requisitos do art. 3º da CLT.

\subsection*{Exemplo 2}

Uma pessoa é chamada uma única vez para reparar um equipamento específico e não se integra à atividade da empresa. O pagamento, sozinho, cria vínculo?

\textbf{Não.} O pagamento indica onerosidade, mas pode faltar não eventualidade e subordinação. O caso se aproxima de prestação eventual ou autônoma.

\section*{Quadro-resumo}

\begin{tabularx}{\linewidth}{lX}
\toprule
 \textbf{Requisito} & \textbf{Pergunta de prova} \\
\midrule
 pessoa física & quem presta é pessoa natural? \\
 pessoalidade & pode se substituir livremente? \\
 não eventualidade & existe continuidade e integração à atividade? \\
 onerosidade & há contraprestação pelo trabalho? \\
 subordinação & o tomador dirige e fiscaliza a prestação? \\
 primazia da realidade & os fatos correspondem ao contrato escrito? \\
\bottomrule
\end{tabularx}


\section*{Questões de fixação}

\subsection*{Questão 1 — (questão inédita)}

É requisito clássico da relação de emprego:

A) autonomia plena do trabalhador. \\ B) eventualidade. \\ C) pessoalidade. \\ D) gratuidade. \\ E) ausência de direção do tomador.

\begin{alerta}
\textbf{Gabarito: C.} Pessoalidade é um dos requisitos do art. 3º da CLT. Autonomia, eventualidade e gratuidade apontam para relações distintas, conforme o caso.
\end{alerta}

\subsection*{Questão 2 — (questão inédita)}

O trabalho realizado em domicílio ou a distância:

A) nunca pode configurar emprego. \\ B) só configura emprego se houver controle presencial. \\ C) pode configurar emprego se presentes os pressupostos da relação. \\ D) é sempre trabalho autônomo. \\ E) elimina a subordinação jurídica.

\begin{alerta}
\textbf{Gabarito: C.} O art. 6º da CLT não distingue o local de trabalho quando os requisitos do vínculo estão presentes.
\end{alerta}

\subsection*{Questão 3 — (questão inédita)}

O princípio da primazia da realidade indica que:

A) o nome dado ao contrato sempre prevalece. \\ B) os fatos efetivamente praticados são relevantes para qualificar a relação. \\ C) todo contrato civil é fraudulento. \\ D) apenas documentos escritos podem provar o vínculo. \\ E) a autonomia é presumida mesmo quando há ordens diárias.

\begin{alerta}
\textbf{Gabarito: B.} A realidade da prestação pode revelar vínculo diverso daquele descrito formalmente.
\end{alerta}

\subsection*{Questão 4 — (questão inédita)}

Subordinação jurídica corresponde:

A) à simples necessidade financeira do trabalhador. \\ B) à obrigação de o trabalhador possuir apenas um cliente. \\ C) ao poder de direção, fiscalização e organização do empregador. \\ D) à inexistência de salário. \\ E) à prestação necessariamente presencial.

\begin{alerta}
\textbf{Gabarito: C.} A subordinação jurídica decorre da sujeição ao poder diretivo do empregador; não se confunde com dependência econômica.
\end{alerta}

\subsection*{Questão 5 — (questão inédita)}

O empregador, conforme a CLT, caracteriza-se por:

A) não assumir nenhum risco da atividade econômica. \\ B) admitir, assalariar e dirigir a prestação pessoal de serviço, assumindo os riscos da atividade. \\ C) não poder fiscalizar a prestação. \\ D) contratar somente pessoas jurídicas. \\ E) atuar sempre sem finalidade econômica.

\begin{alerta}
\textbf{Gabarito: B.} É a definição central do art. 2º da CLT.
\end{alerta}

\section*{Revisão final}

\begin{itemize}
  \item[$\square$] Memorizei os cinco requisitos do vínculo.
  \item[$\square$] Sei explicar pessoalidade e não eventualidade.
  \item[$\square$] Diferencio subordinação jurídica de dependência econômica.
  \item[$\square$] Sei que trabalho remoto também pode ser emprego.
  \item[$\square$] Conheço a primazia da realidade.
  \item[$\square$] Sei que o empregador assume os riscos da atividade.
  \item[$\square$] Diferencio empregado, autônomo e eventual.
  \item[$\square$] Sei que estágio irregular pode ser desvirtuado.
\end{itemize}

\end{document}
